\chapter{Presupuesto}

Para la elaboración del presupuesto se tuvieron en cuenta los costos directos de adquisición de componentes electrónicos y servomotores, así como los costos indirectos prorrateados de uso de infraestructura institucional como osciloscopios, fuentes de alimentación, estaciones de soldadura, computadores y la plataforma robótica cuadrúpedo. Los precios se estimaron con referencias de tiendas de electrónica y plataformas de comercio electrónico en Colombia, considerando valores promedio para componentes de pequeña escala (resistencias, headers, conectores) y precios de mercado actuales para elementos más costosos como la Raspberry Pi 5 y servomotores. La infraestructura del laboratorio se considera propiedad institucional, por lo que no se incluye su costo completo, aunque se menciona la posibilidad de un prorrateo institucional según el uso por estudiantes o proyectos.

\section{Costos directos del proyecto}

\begin{table}[H]
\centering
\caption{Listado de materiales y costos del proyecto}
\label{tab:costos_materiales}
\begin{tabular}{|c|l|c|c|}
\hline
\textbf{Ítem} & \textbf{Descripción} & \textbf{Cantidad} & \textbf{Total (COP)} \\
\hline
1  & Raspberry Pi 5 16GB & 1 & 1.035.300 \\
2  & Módulo Sensor BNO055 – IMU 9DOF con MCU Integrado & 1 & 65.000 \\
3  & Sensor de Temperatura Salida Digital DS18B20 & 12 & 5.500 \\
4  & Servomotores & 12 & 384.000 \\
5  & Female Header 2.54 mm 1x14 & 1 & 2.000 \\
6  & Impresión de PCB & 1 & 50.000 \\
7  & Módulo PCA9685 (Controlador PWM) & 1 & 35.000 \\
8  & Regulador DC-DC LM2596ADJ (3.2V–35V) & 1 & 15.000 \\
\hline
\multicolumn{3}{|r|}{\textbf{Total costos directos}} & \textbf{1.591.800 COP} \\
\hline
\end{tabular}
\end{table}


\section{Costos indirectos prorrateados}

En el desarrollo del proyecto se utilizarán equipos del laboratorio de robótica, tales como:

\begin{itemize}
\item Osciloscopio
\item Fuente de alimentación de laboratorio
\item Estación de soldadura
\item Computadores de laboratorio
\item Plataforma robótica cuadrúpeda
\end{itemize}

Estos equipos pertenecen a la infraestructura institucional, por lo que no se incluye su costo total
en el presupuesto. En caso de considerarse un prorrateo institucional, este se calcularía según el
número de estudiantes o proyectos que utilizan dichos equipos durante su vida útil.

Para efectos de este anteproyecto, no se asignan costos prorrateados adicionales.

\section{Total general del proyecto}

\begin{table}[h]
\centering
\caption{Resumen de costos del proyecto}
\label{tab:resumen_costos}
\begin{tabular}{|l|c|}
\hline
\textbf{Categoría} & \textbf{Valor (COP)} \\
\hline
Costos directos & 2.029.270 \\
Costos prorrateados & 0 \\
\hline
\textbf{TOTAL GENERAL} & \textbf{2.029.270 COP} \\
\hline
\end{tabular}
\end{table}


\section{Supuestos y criterios de cálculo}

\begin{itemize}

\item Los precios se estimaron con referencias de tiendas de electrónica en Colombia y
plataformas de comercio electrónico.

\item Los componentes electrónicos pequeños (resistencias, headers, conectores) se
calcularon con precios promedio de kits comerciales.

\item El precio de la Raspberry Pi 5 compatible se estimó alrededor de 1.035.300 COP,
valor común en el mercado colombiano.

\item La plataforma robótica cuadrúpeda y los servomotores corresponden a valores
definidos dentro del laboratorio.

\item Se considera que la infraestructura del laboratorio (equipos de medición,
computadores y herramientas) es propiedad institucional, por lo que no se incluye su
costo completo dentro del presupuesto del proyecto.

\end{itemize}
